\documentclass[11pt,a4paper]{article}
\usepackage[utf8]{inputenc}\usepackage[T1]{fontenc}\usepackage[italian,english]{babel}
\usepackage{amsmath,amssymb,amsthm}\usepackage[margin=2.2cm]{geometry}\usepackage{longtable,booktabs,array,calc}
\usepackage{listings,xcolor}\usepackage{hyperref}\usepackage{url}
\providecommand{\tightlist}{\setlength{\itemsep}{0pt}\setlength{\parskip}{0pt}}
\lstset{basicstyle=\ttfamily\scriptsize,breaklines=true,frame=single,captionpos=b,columns=fullflexible,keepspaces=true,inputencoding=utf8,extendedchars=true}
\title{Proof Pursuit --- Problem 1, part 1\\ \large prove, decisioni degli agenti, codice verificato, letteratura}
\author{Team Proof Pursuit (Researcher: T. Tumini; Referee: gabundos; Creative: M. Renzi; gio3r3gio)}
\date{\today}
\begin{document}\maketitle\tableofcontents
\section*{Come leggere questo rapporto}
Per ogni cella: la richiesta ufficiale, lo stato dichiarato dal team (mai ``risolto'' senza revisione umana), la consegna
con la prova, poi la traccia delle decisioni degli agenti (Researcher: famiglia e motivo; Referee: verdetto ed errore
fatale; Creative: quando e perché è intervenuto) e il codice su cui il tentativo poggia con l'esito della riesecuzione
fidata. DIMOSTRATO, CITATO, VERIFICATO SU INTERVALLO FINITO e CONGETTURA restano distinti. L'architettura del sistema
è descritta nel README del repository.
\section{Problem 1 — Angles between lines (Fejes Toth)}
\subsection{Cella 1 (1 punti) — stato: revisionato}
\subparagraph{Parte 1 --- Lines in the
plane}\label{parte-1-lines-in-the-plane}

\textbf{Punteggio:} 1 point · \textbf{Valutazione:} Judged

We begin in the plane (\(d = 2\)), where the conjectured optimum splits
the \(N\) lines as evenly as possible between two perpendicular
directions. Let \(\ell_1, \ldots, \ell_N\) be lines in \(\mathbb{R}^2\).
Prove that \[
S(\ell_1, \ldots, \ell_N) \;\le\; \frac{\pi}{2} \cdot \left\lfloor \frac{N^2}{4} \right\rfloor .
\]

\textbf{Enunciato protetto dato agli agenti (state.json).} \texttt{Let l\_1..l\_N be lines through the origin in R\textasciicircum{}2. Prove S(l\_1..l\_N) = sum\_\{i\textless{}j\} theta(l\_i,l\_j) \textless{}= (pi/2) * floor(N\textasciicircum{}2/4).}

\textbf{Evidenza registrata in STATUS.md.} prova completa (semigiro casuale, argomento E) scritta a mano; Referee READY\_FOR\_HUMAN (A PASS, B PASS); approvata; submission/parte-1.md

\subsubsection{Consegna (prova)}
\paragraph{Problema 1 --- Parte 1 --- bozza di
consegna}\label{problema-1-parte-1-bozza-di-consegna}

\textbf{Stato dichiarato: RISOLTA (approvata).} Prova completa scritta
dal team; Referee automatico READY\_FOR\_HUMAN (giudice A matematica
PASS, giudice B evidenze PASS); approvazione umana registrata in
\texttt{runs/p1\_c1/approval.json}.

\subparagraph{1. Risultato e ambito}\label{risultato-e-ambito}

Somma degli angoli fra rette nel piano: l'enunciato della parte è
dimostrato integralmente, nella generalità richiesta (vedi la prova per
le ipotesi esatte). Nessun calcolo al computer è necessario.

\subparagraph{2. Dimostrazione}\label{dimostrazione}

\paragraph{\texorpdfstring{Cell 1: the sum of the angles between \(N\)
lines in the
plane}{Cell 1: the sum of the angles between N lines in the plane}}\label{cell-1-the-sum-of-the-angles-between-n-lines-in-the-plane}

\textbf{Target.} For lines \(\ell_1,\dots,\ell_N\) through the origin of
\(\mathbb R^2\) (repetitions allowed),
\(S=\sum_{i<j}\theta(\ell_i,\ell_j)\le \frac{\pi}{2}\left\lfloor \frac{N^2}{4}\right\rfloor\),
where \(\theta(\ell,\ell')=\arccos|\langle x,x'\rangle|\) for unit
vectors \(x,x'\) spanning \(\ell,\ell'\).

\subparagraph{1. Directions and angles}\label{directions-and-angles}

Every line \(\ell\) through the origin of \(\mathbb R^2\) is spanned by
exactly one unit vector of the form \((\cos t,\sin t)\) with
\(t\in[0,\pi)\) (the two unit vectors of \(\ell\) are
\(\pm(\cos t,\sin t)\), and exactly one of the two angles \(t,t+\pi\)
lies in \([0,\pi)\)). Call \(t=\tau(\ell)\) the \emph{direction} of
\(\ell\). For two lines with directions \(t,t'\in[0,\pi)\) we have
\(\langle(\cos t,\sin t),(\cos t',\sin t')\rangle=\cos(t-t')\), hence
\(\theta(\ell,\ell')=\arccos|\cos(t-t')|\). Put \(u=|t-t'|\in[0,\pi)\).
If \(u\le\pi/2\) then \(\cos u\ge 0\), so \(|\cos u|=\cos u\) and
\(\arccos(\cos u)=u\) because \(u\in[0,\pi]\). If \(u>\pi/2\) then
\(\cos u<0\), so \(|\cos u|=-\cos u=\cos(\pi-u)\) with
\(\pi-u\in(0,\pi/2)\), and \(\arccos(\cos(\pi-u))=\pi-u\). In both cases
\[\theta(\ell,\ell')=\min(u,\pi-u),\qquad u=|\tau(\ell)-\tau(\ell')|.\tag{1}\]
Coincident lines have \(u=0\) and angle \(0\), as they should. Note also
that \(\min(u,\pi-u)\) is unchanged if \(u\) is replaced by \(\pi-u\),
so (1) also holds with \(u=(\tau(\ell')-\tau(\ell))\bmod\pi\) (the
representative in \([0,\pi)\)): indeed this quantity equals \(|t-t'|\)
or \(\pi-|t-t'|\).

\subparagraph{2. A random half-turn of
directions}\label{a-random-half-turn-of-directions}

For \(x\in\mathbb R\) write \(x\bmod\pi\) for the unique representative
of \(x+\pi\mathbb Z\) in \([0,\pi)\). For \(\psi\in[0,\pi)\) define
\[A_\psi=\{t\in[0,\pi):\ (t-\psi)\bmod\pi\in[0,\pi/2)\}.\] Let \(\psi\)
be a random variable with the uniform distribution on \([0,\pi)\),
i.e.~\(\Pr[\psi\in B]=|B|/\pi\) for every Borel set
\(B\subseteq[0,\pi)\), where \(|B|\) is Lebesgue measure.

\textbf{Lemma 1.} For fixed \(t,t'\in[0,\pi)\) let \(u=(t'-t)\bmod\pi\)
and \(\varphi=\min(u,\pi-u)\). Then
\[\Pr\big[\text{exactly one of } t,t' \text{ lies in } A_\psi\big]=\frac{2\varphi}{\pi}.\]

\emph{Proof.} Step 1 (reduction to \(t=0\)). Put
\(\psi'=(\psi-t)\bmod\pi\). The map \(x\mapsto(x-t)\bmod\pi\) is a
bijection of \([0,\pi)\) onto itself which is a translation on each of
the two pieces \([0,t)\) and \([t,\pi)\), hence it preserves Lebesgue
measure; therefore \(\psi'\) is again uniform on \([0,\pi)\). Moreover
\((t-\psi)\bmod\pi=(-(\psi-t))\bmod\pi=(-\psi')\bmod\pi=(0-\psi')\bmod\pi\)
and \((t'-\psi)\bmod\pi=((t'-t)-(\psi-t))\bmod\pi=(u-\psi')\bmod\pi\).
Hence \(t\in A_\psi\iff 0\in A_{\psi'}\) and
\(t'\in A_\psi\iff u\in A_{\psi'}\). So it suffices to prove the claim
for the pair \((0,u)\) with \(u\in[0,\pi)\) and a uniform \(\psi\) (we
rename \(\psi'\) as \(\psi\)).

Step 2 (the two sets). For \(s\in[0,\pi)\) let
\(J_s=\{\psi\in[0,\pi): s\in A_\psi\}=\{\psi: (s-\psi)\bmod\pi<\pi/2\}\).
Since \(s-\psi\in(-\pi,\pi)\): if \(\psi\le s\) then
\((s-\psi)\bmod\pi=s-\psi\) and the condition reads \(\psi>s-\pi/2\); if
\(\psi>s\) then \((s-\psi)\bmod\pi=s-\psi+\pi\) and the condition reads
\(\psi>s+\pi/2\). Therefore
\[J_s=\big(\max(0,s-\tfrac{\pi}{2}),\,s\big]\ \cup\ \big(s+\tfrac{\pi}{2},\,\pi\big)\cap[0,\pi),\]
where the second interval is empty when \(s+\pi/2\ge\pi\). In particular
\(|J_s|=\pi/2\) for every \(s\): if \(s<\pi/2\) the two pieces have
lengths \(s\) and \(\pi/2-s\); if \(s\ge\pi/2\) the first piece has
length \(\pi/2\) and the second is empty. For \(s=0\) the first piece is
\(\{0\}\) (measure \(0\)) and \(J_0=\{0\}\cup(\pi/2,\pi)\).

Step 3 (the symmetric difference). The event ``exactly one of \(0,u\)
lies in \(A_\psi\)'' is \(\{\psi\in J_0\triangle J_u\}\), and
\(|J_0\triangle J_u|=|J_0|+|J_u|-2|J_0\cap J_u|=\pi-2|J_0\cap J_u|\).
\emph{Case \(u\le\pi/2\).} Then \(J_u=(0,u]\cup(u+\pi/2,\pi)\) (the
second piece is empty if \(u=\pi/2\)). Up to the null set \(\{0\}\),
\(J_0=(\pi/2,\pi)\); since \(u\le\pi/2\),
\((0,u]\cap(\pi/2,\pi)=\emptyset\), so \(J_0\cap J_u=(u+\pi/2,\pi)\) up
to a null set, of measure \(\pi/2-u\). Hence
\(|J_0\triangle J_u|=\pi-2(\pi/2-u)=2u=2\varphi\). \emph{Case
\(u>\pi/2\).} Then \(u+\pi/2>\pi\), so \(J_u=(u-\pi/2,u]\) with
\(0<u-\pi/2<\pi/2\). Thus
\(J_0\cap J_u=(\pi/2,\pi)\cap(u-\pi/2,u]=(\pi/2,u]\) up to a null set,
of measure \(u-\pi/2\). Hence
\(|J_0\triangle J_u|=\pi-2(u-\pi/2)=2(\pi-u)=2\varphi\). In both cases
the probability is \(|J_0\triangle J_u|/\pi=2\varphi/\pi\). \(\square\)

\subparagraph{\texorpdfstring{3. Counting the pairs separated by
\(A_\psi\)}{3. Counting the pairs separated by A\_\textbackslash psi}}\label{counting-the-pairs-separated-by-a_psi}

Let \(t_i=\tau(\ell_i)\) and, for \(\psi\in[0,\pi)\),
\(k_\psi=\#\{i: t_i\in A_\psi\}\in\{0,1,\dots,N\}\). A pair \(\{i,j\}\)
has exactly one member with direction in \(A_\psi\) if and only if one
index is among the \(k_\psi\) ``inside'' indices and the other among the
\(N-k_\psi\) ``outside'' indices; hence the number of such pairs is
exactly \(k_\psi(N-k_\psi)\). Writing \(X_{ij}(\psi)\) for the indicator
of the event of Lemma 1 for the pair \((t_i,t_j)\), we have
\(k_\psi(N-k_\psi)=\sum_{i<j}X_{ij}(\psi)\) for every \(\psi\) (each
\(X_{ij}\) is the indicator of a finite union of intervals, so
everything is measurable), and by linearity of expectation, Lemma 1 and
(1),
\[\mathbb E\big[k_\psi(N-k_\psi)\big]=\sum_{i<j}\Pr[X_{ij}=1]=\sum_{i<j}\frac{2}{\pi}\theta(\ell_i,\ell_j)=\frac{2}{\pi}S.\tag{2}\]

\subparagraph{4. Conclusion}\label{conclusion}

For every integer \(k\) we have \(N^2-4k(N-k)=(N-2k)^2\ge0\), so
\(k(N-k)\le N^2/4\); as \(k(N-k)\) is an integer,
\(k(N-k)\le\lfloor N^2/4\rfloor\). Applying this to \(k=k_\psi\) for
every \(\psi\) and taking expectations, (2) gives
\[S=\frac{\pi}{2}\,\mathbb E\big[k_\psi(N-k_\psi)\big]\le\frac{\pi}{2}\left\lfloor\frac{N^2}{4}\right\rfloor.\qquad\blacksquare\]

\emph{Remark (not needed for the cell).} The bound is attained: put
\(\lceil N/2\rceil\) lines along the \(x\)-axis and
\(\lfloor N/2\rfloor\) along the \(y\)-axis; the only nonzero angles are
the \(\lceil N/2\rceil\lfloor N/2\rfloor=\lfloor N^2/4\rfloor\) cross
pairs, each equal to \(\pi/2\).

\subparagraph{3. Verifica: istruzioni, dipendenze,
tempi}\label{verifica-istruzioni-dipendenze-tempi}

Prova a mano, nessuna dipendenza. Verifica meccanica non necessaria; il
Referee automatico (due giudici indipendenti,
\texttt{runs/p1\_c1/attempts/packet\_001.json}) non ha trovato passaggi
non giustificati.

\subparagraph{4. Fonti e contributo}\label{fonti-e-contributo}

\begin{itemize}
\tightlist
\item
  arXiv:1801.07837 Bilyk, Matzke, On the Fejes Toth problem about the
  sum of angles between lines (context only: the planar case is the only
  settled case; the proof below is self-contained and does not use the
  paper)
\item
  Posizione rispetto allo stato dell'arte: Bilyk-Matzke {[}1801.07837{]}
  state that d=1 (the plane) is the only settled case and give several
  proofs, one of which is a Stolarsky-type identity; our argument is of
  the same nature but written independently and in full, so nothing is
  cited as a black box. Lim-McCann {[}2007.08698{]} concern a
  one-parameter deformation and are not used.
\item
  Contributo: la prova qui scritta è del team, per esteso.
\end{itemize}

\subparagraph{5. Limiti e parti
irrisolte}\label{limiti-e-parti-irrisolte}

Nessuna: la parte è dimostrata. Gap dichiarati: nessuno.

\subsubsection{Decisioni degli agenti}
\paragraph{Run \texttt{p1\_c1}}
\begin{description}
\item[Iterazione 1 — attempt\_001]
\textbf{Researcher}: famiglia \texttt{direct\_proof}, sotto-obiettivo: Prove the planar bound by a random-half-turn (Crofton/Stolarsky-type) double counting: S = (pi/2) E[k(N-k)].. Stato dichiarato: \texttt{CELL\_SOLVED\_CANDIDATE}.
\emph{Perché questa scelta}: The planar case admits an elementary integral-geometric proof: the angle between two lines is proportional to the probability that a random half-turn of directions separates them, and a separator splits N lines into k and N-k, so the sum is (pi/2) E[k(N-k)] \textless{}= (pi/2) floor(N\textasciicircum{}2/4). Chosen because it is short, exact, and needs no case analysis on the configuration.
\emph{Rispetto allo stato dell'arte}: Bilyk-Matzke [1801.07837] state that d=1 (the plane) is the only settled case and give several proofs, one of which is a Stolarsky-type identity; our argument is of the same nature but written independently and in full, so nothing is cited as a black box. Lim-McCann [2007.08698] concern a one-parameter deformation and are not used.
\textbf{Referee}: verdetto \texttt{UNKNOWN\_STATUS} (review\_status \texttt{READY\_FOR\_HUMAN}).
\emph{Prossimo passo richiesto}: Human reviews the exact target, proof and evidence, then approves explicit claims
\end{description}

\subsubsection{Codice su cui poggia il tentativo}
Nessun codice nei tentativi.

\section{arXiv literature consulted}
\begin{thebibliography}{99}
\bibitem{1801.07837v1} Dmitriy Bilyk, Ryan W Matzke. \emph{On the Fejes Tóth Problem about the Sum of Angles Between Lines}. arXiv:1801.07837v1 (2018). \url{http://arxiv.org/abs/1801.07837v1}
\bibitem{2007.08698v2} Tongseok Lim, Robert J. McCann. \emph{On Fejes Tóth's conjectured maximizer for the sum of angles between lines}. arXiv:2007.08698v2 (2020). \url{http://arxiv.org/abs/2007.08698v2}
\bibitem{1204.3850v1} Jérémie Chalopin, Shantanu Das, Yann Disser, Matúš Mihalák, Peter Widmayer. \emph{Simple Agents Learn to Find Their Way: An Introduction on Mapping Polygons}. arXiv:1204.3850v1 (2012). \url{http://arxiv.org/abs/1204.3850v1}
\end{thebibliography}
\end{document}
